% Source: midterm3_final_v1.html
\documentclass[11pt]{article}
\usepackage[utf8]{inputenc}
\usepackage[T1]{fontenc}
\usepackage[margin=1in]{geometry}
\usepackage{amsmath}
\usepackage{amssymb}
\usepackage{siunitx}
\usepackage{enumitem}
\usepackage{parskip}

\title{PS160 -- midterm3\_final\_v1.html}
\date{}

\begin{document}
\maketitle

\section*{Questions}

\begin{enumerate}[leftmargin=*]
  \item A horizontal pipe of radius 0.24 m carries water at 2.5 m/s and pressure of 227 kPa. If the pipe narrows to a radius of 0.15 m, calculate the pressure \emph{p} in kPa in the narrow pipe. The density of water is 1000 kg/m\textsuperscript{3}.

\emph{p} = kPa \hfill \textit{(10 pts, Numeric)}

  \item A simple harmonic oscillator consists of a block of an unknown mass in kilograms and spring with spring constant 660 N/m on a frictionless horizontal surface. The block is moved a distance 1.48 m\textbf{}from equilibrium and released from rest. If the block's maximum speed at the equilibrium point is 9.4 m/s, calculate the mass of the block in kg. \hfill \textit{(10 pts, Numeric)}

  \item A 7-kg block attached to a spring of force constant 132 N/m undergoes simple harmonic motion. At a particular time the mass' speed is 2.12 m/s and its position is 0.26 m. Calculate the amplitude of its oscillation.

\emph{A}= m \hfill \textit{(10 pts, Numeric)}

  \item Suppose a wave is given by \emph{y}(\emph{x},\emph{t})=5 cos(0.2$\pi$\emph{t -}0.3$\pi$\emph{x} + 0.25$\pi$) with all quantities in MKS units. Determine the following in MKS units:

\textbf{(a)} amplitude [DROPDOWN: 6 m | 2.5 m | 5 m | 2 m | 3 m]
;

\textbf{(b)} phase shift [DROPDOWN: 0.785 rad | 0.628 rad | 0.125 rad | 0.25 rad]
;

\textbf{(c)} wave number [DROPDOWN: 0.628 rad/m | 0.942 rad/m | 0.785 rad/m | 0.15 rad/m]
;

\textbf{(d)} angular frequency [DROPDOWN: 0.628 rad/s | 0.314 rad/s | 0.2 rad/s | 0.942 rad/s]
;

\textbf{(e)} frequency [DROPDOWN: 0.15 Hz | 0.1 Hz | 5 Hz | 0.125 Hz]
;

\textbf{(f)} period [DROPDOWN: 0.942/m | 20 s | 5 s | 0.628 s | 0.628/m | 0.125/m | 6.67 s | 0.15 /m | 10 s]
;

\textbf{(g)} wavelength [DROPDOWN: 5 m | 6.67 m | 0.942 m | 10 m]

\textbf{(h)} propagation speed [DROPDOWN: 3.14 m/s | 1.50 m/s | 2.35 m/s | 0.667 m/s]
;\textbf{}

\textbf{(i)}transverse \textbf{speed} at \emph{x} = 0, \emph{t} = 0: [DROPDOWN: 2.22 m/s | 10 m/s | 3.14 m/s | 6.67 m/s] \hfill \textit{(10 pts, Dropdown)}

  \item A pipe closed at one end has a length of 1 m. If the temperature of the room is 8$^{\circ}$C, calculate the fundamental frequency in Hz.

Take the speed of sound to be $v=(331\;\mathrm{m/s})\sqrt{\frac{T}{273\;\mathrm{K}}}$ . \\ \hfill \textit{(10 pts, Numeric)}

  \item A motorcyclist is traveling 21 m/s due west relative the ground. A fire truck is traveling on the same road, going east at 39 m/s and away from the motorcyclist, having already passed it going the opposite direction. If the fire truck siren sounds at 1,071 Hz, what frequency does the motorcyclist hear? Sound speed in air: 340 m/s. \hfill \textit{(10 pts, Numeric)}

  \item A metal has a coefficient of linear expansion of 1.2 $\times$ 10\textsuperscript{-5} (C$^{\circ}$)\textsuperscript{-1}. Suppose a pipe made of this metal is 5.3 m long. Calculate its length increase in mm if it is warmed by 77 C$^{\circ}$. \\

$\Delta$\emph{L}= mm \\ \hfill \textit{(10 pts, Numeric)}

  \item Suppose a block of a substance has mass 4.7 kg and is in solid form at its melting point of 3$^{\circ}$C. How much energy in kJ must be transferred to the block in order to put it in its liquid form at the same temperature? The latent heat of fusion of the substance is \emph{L}\textsubscript{f} = 117,005 J/kg.

\emph{Q} = kJ \hfill \textit{(10 pts, Numeric)}

  \item An ideal gas has an initial pressure of \emph{p}\textsubscript{0} and volume of 2.2 m\textsuperscript{3}. It undergoes a process whereby the pressure changes to (2.4)\emph{p}\textsubscript{0} while the Kelvin temperature increases by 58\%. Calculate the final volume in meters cubed. \hfill \textit{(10 pts, Numeric)}

  \item In the given \emph{pV}-diagram (not to scale), \emph{p}\textsubscript{1} = 106.8 kPa, \emph{p}\textsubscript{2} = 223.9 kPa; \emph{V}\textsubscript{1} = 0.9 m\textsuperscript{3}, \emph{V}\textsubscript{2} = 1.8 m\textsuperscript{3}, and \emph{V}\textsubscript{3} = 3.5 m\textsuperscript{3}. Calculate the work in kJ done \textbf{on} the system during these processes.

[Image: ME19\_PV\_3.png] \\

\emph{W}\textsubscript{on}\emph{}= kJ \hfill \textit{(10 pts, Numeric)}

  \item A diatomic ideal gas at constant volume of 1.5 m\textsuperscript{3} decreases pressure from 136 kPa to 81 kPa. Calculate the energy transfer by heat in kJ during this process. \\ \hfill \textit{(10 pts, Numeric)}

  \item A monatomic ideal gas expands adiabatically from 1.00 m\textsuperscript{3} to 2.7 m\textsuperscript{3}. If the initial pressure is 244 kPa, calculate the work done \textbf{by} the system in kJ. \hfill \textit{(10 pts, Numeric)}

  \item An engine with an efficiency 0.2 performs 1,041 J of work during a typical cycle. How much energy does it exhaust to the cold reservoir in one cycle? \hfill \textit{(10 pts, Numeric)}

  \item A beam of light is incident on a mirror lying flat on the ground at an angle of $\theta$ = 70$^{\circ}$ with respect to the normal. A second mirror is aligned with one edge of the first mirror and tilted up at an angle of $\phi$ = 59$^{\circ}$. At what angle $\alpha$ with respect to the normal does the ray reflect off the second mirror?

[Image: ME33\_plane\_mirror\_1.png]

$\alpha$ = $^{\circ}$ \hfill \textit{(10 pts, Numeric)}

  \item A ray of light traveling in air is incident at an angle of 27$^{\circ}$ with respect to the normal on a fluid with index of refraction 1.7. Calculate the angle of refraction in degrees. \hfill \textit{(10 pts, Numeric)}

  \item A liquid-air boundary has a critical angle of 43.4$^{\circ}$ for total internal reflection. If a ray of light traveling in air has an incident angle of 22$^{\circ}$ at the liquid-air boundary, calculate the angle between the refracted ray and the surface normal. \hfill \textit{(10 pts, Numeric)}

  \item Unpolarized light of intensity 20 W/cm\textsuperscript{2} is incident on a set of three polarizing filters, rotated 24$^{\circ}$, 49$^{\circ}$, and 24$^{\circ}$ from the vertical, respectively. Calculate the light intensity in W/cm\textsuperscript{2} leaving the third filter. \hfill \textit{(10 pts, Numeric)}

  \item An object is located at 8.00 cm in front of a concave spherical mirror of focal length 14.0 cm. The image position is [DROPDOWN: -18.7 cm | 18.7 cm | 12.4 cm | 5.09 cm | -12.4 cm | -9.8 cm | 9.8 cm | -5.09 cm]
cm, the magnification is [DROPDOWN: 2.33 | 4.66 | -2.33 | -3.28 | -4.66 | 3.28]
, the image is [DROPDOWN: real | more information required | virtual]
and [DROPDOWN: uright | more information required | inverted]
. \hfill \textit{(10 pts, Dropdown)}

  \item A thin converging lens has a focal length of 12.0 cm. If the object is located at 9.00 cm, the image position is [DROPDOWN: -36 | 36 | 34 | -24 | -18]
cm, the magnification is [DROPDOWN: -6.3 | 3 | 6.3 | 4 | -3 | 0.333]
, and the image is both [DROPDOWN: real | virtual]
and [DROPDOWN: inverted | upright]
. \hfill \textit{(10 pts, Dropdown)}

  \item A thin diverging lens has a focal length of 15.0 cm. If the object distance is 10.0 cm, the image distance is [DROPDOWN: -0.167 | -7.86 | -6 | 6 | -4.82]
cm, the magnification is [DROPDOWN: -0.9 | 0.6 | -0.6 | 0.9 | 0.72]
, and the image is both [DROPDOWN: real | virtual | something else]
and [DROPDOWN: upright | inverted | both upright and inverted]
. \hfill \textit{(10 pts, Dropdown)}

  \item The figure shows a contact lens made from a material with an index of refraction of 1.5. The front and back surfaces have radii of curvature of 2.2 cm and 2.7 cm, respectively. Calculate the lens' focal length in cm.

[Image: contact\_lens.png] \hfill \textit{(10 pts, Numeric)}

  \item A silicon solar cell with an index of refraction of 3.88 is coated with an anti-reflection thin film having an index of refraction of 1.45. Calculate the minimum film thickness in nm that produces the least reflection at a wavelength of 557 nm. \hfill \textit{(10 pts, Numeric)}

  \item A single slit experiment involves monochromatic, coherent light of wavelength 639 nm incident on an aperture of width 0.23 mm. If the screen is 1.7 m away, determine the distance in mm from the central maximum to the \emph{m}= 4 dark line. \hfill \textit{(10 pts, Numeric)}

  \item It is found that monochromatic, coherent light of wavelength 628 nm incident on a pair of slits creates a central maximum with width 9.3 mm. If the screen is 1.5 m away, determine the separation \emph{d} of the two slits in mm.

[The width of the central maximum is defined as the distance between the \emph{m}=0 dark bands on either side of the \emph{m}=0 maximum.] \hfill \textit{(10 pts, Numeric)}

  \item At what grazing angle in degrees from the surface will x-rays of wavelength 0.28 nm on a crystal with interatomic spacings between the Bragg planes of 0.55 nm result in a first-order interference fringe? \\ \hfill \textit{(10 pts, Numeric)}

\end{enumerate}

\end{document}